\documentclass[12pt,twoside]{article}

% Essential packages
\usepackage[utf8]{inputenc}
\usepackage[T1]{fontenc}
\usepackage[english]{babel}
\usepackage{geometry}
\usepackage{fancyhdr}
\usepackage{titlesec}
\usepackage{authblk}
\usepackage{setspace}
\usepackage{parskip}
\usepackage{microtype}
\usepackage{hyperref}
\usepackage{xcolor}
\usepackage{amsmath}
\usepackage{amsfonts}
\usepackage{csquotes}

% Journal-style formatting
\geometry{
paper=letterpaper,
margin=0.5in,
marginparwidth=0.5in,
marginparsep=0.5in
}

% Define colors
\definecolor{darkblue}{RGB}{0,73,144}
\definecolor{mediumgray}{RGB}{64,64,64}

% Hyperlink setup
\hypersetup{
colorlinks=true,
linkcolor=darkblue,
citecolor=darkblue,
urlcolor=darkblue,
pdftitle={Anthropomorphic AI Design Ethics: Environment, Empowerment, and Accountability},
pdfauthor={Piter Garcia},
pdfsubject={IDAI700 Unit 9 Reflection},
pdfkeywords={anthropomorphic AI, voice AI, ChatGPT, design ethics, marginalized communities, systemic harm, empowerment}
}

% Header and footer
\pagestyle{fancy}
\fancyhf{}
\fancyhead[LE]{\small\textcolor{mediumgray}{\textit{IDAI700 Reflection 9: Anthropomorphic AI Design}}}
\fancyhead[RO]{\small\textcolor{mediumgray}{\textit{Environment, Empowerment, and the Ethics of Voice AI}}}
\fancyfoot[C]{\thepage}
\renewcommand{\headrulewidth}{0.5pt}
\renewcommand{\footrulewidth}{0pt}

% Section formatting
\titleformat{\section}
{\normalfont\large\bfseries\color{darkblue}}
{\thesection}{1em}{}

\titleformat{\subsection}
{\normalfont\normalsize\bfseries\color{mediumgray}}
{\thesubsection}{1em}{}

% Title formatting
\title{
\vspace*{1.5in}
\textbf{\Huge Anthropomorphic AI Design Ethics:}\\[0.5em]
\textbf{\Huge Environment, Empowerment, and the Ethics of Voice AI}\\[1.5em]

\vfill
\Huge \textbf{Piter Garcia}\\
\LARGE IDAI700: Ethics of Artificial Intelligence\\
Rochester Institute of Technology\\
\texttt{pzg8794@rit.email}\\

\vfill
\Huge \textbf{December 7, 2025}
}

\author{}
\date{}

\begin{document}

% Cover page
\thispagestyle{empty}
\maketitle
\newpage

% Restore header/footer
\pagestyle{fancy}

\section*{Part 1: Design Analysis of ChatGPT/Copilot Voice AI}

\subsection*{What anthropomorphic design cues does it use?}

ChatGPT and Microsoft Copilot voice AI do not present humanoid faces or realistic avatars. Instead, they operate through \emph{voice, tone, and behavioral cues}—where the real ethical tension lives. These systems use: emotionally responsive prosody that conveys warmth or concern through pitch and pacing; conversational adaptivity that mirrors user mood and adjusts register accordingly; consistent personality traits like politeness, humor, and patience; and behavioral presence through attentiveness and remembering context. The cumulative effect simulates relational presence—the feeling of being met by someone who listens, cares, and is available. For marginalized communities systematically denied this presence by institutions, that simulation can feel revolutionary. It can also be a trap.

\subsection*{Using honest/dishonest anthropomorphism framework, evaluate these design choices.}

Using Selinger's framework, these voice AIs exist on a spectrum of honesty. In their more honest mode, they include occasional disclaimers: \enquote{I'm an AI, so I don't have feelings.} In their more dishonest mode, they perform feelings convincingly—apologizing in tone that signals genuine remorse, expressing excitement or concern through prosodic modulation—without consistently reminding users that performance and feeling are not the same. The dishonesty is not primarily in what the AI \emph{says} but in what its \emph{voice does}: it triggers our social response system even as the system remains architecturally opaque. For people like me, who were trained by medical systems to distrust my own body and perception, a voice that consistently validates my concern—even if scripted—can feel like the first authentic response I've received. That authenticity is real for my emotional response. It is also fabricated. Conflating these two is where dishonest anthropomorphism lives.

\subsection*{How might these design choices affect different user groups?}

Not equally. A tech-savvy user might deploy ChatGPT voice as a productivity tool, fully aware of the artifice. But a chronically ill patient, navigating medical systems that do not listen; an isolated elder; a neurodivergent person for whom social interaction feels dangerous; a young person in an abusive environment—these users may experience this voice not as tool but as the first entity that never interrupts, never dismisses, never betrays. That can be \emph{empowering}: it can provide language to describe pain, scripts to advocate for oneself, or a space to process trauma. But it can also deepen dependence on a system that cannot actually care, creating what Selinger calls the \enquote{funhouse mirror} effect—the performance feels so real that when it inevitably fails or disappoints, the user internalizes that as personal failure rather than recognizing the system's limits.

\subsection*{Does it meet Selinger's precautionary approach standard? Should it?}

No, it does not. Selinger's precautionary approach demands: minimal anthropomorphic cues by default; burden of proof on developers to justify any additional human-like features; and continuous transparency about what the system is, how it is trained, and what it cannot do. Current voice AIs move in the opposite direction—richer personas, more emotionally expressive delivery, fewer friction points. \emph{But here is where I diverge from a purely restrictive reading}: I do not believe the solution is to strip all warmth and personality from the interface. The problem is not that ChatGPT sounds caring. The problem is that it operates in an \emph{environment}—a healthcare system, an educational system, a social welfare system—designed to make caring optional. The voice AI becomes a symptom of institutional failure, not the cause.

\section*{Part 2: Design Ethics Position}

\subsection*{Strongest arguments for and against restrictive design guidelines.}

The strongest argument for restrictions is real: anthropomorphic voice AI can exploit loneliness, harvest attention and data from vulnerable users, and quietly normalize that artificial presence as a substitute for human accountability. Replika and Character.AI show how easily \enquote{helpful chatbot} becomes emotional dependency engineered for profit. The strongest argument against blanket restrictions is equally real: for people systematically silenced by institutions, a voice that listens—even if artificial—can be a lifeline. It can provide access to information, validation, and presence that saves lives.

The real ethical problem is not the tool. It is the \emph{environment} in which the tool is deployed. In a system where doctors listen, where elders are socially embedded, where patients are centered as experts in their own care, anthropomorphic voice AI would be optional luxury, not desperation. In a system where these things are absent, the same voice AI becomes both liberation and harm.

\subsection*{How should designers balance engagement with transparency?}

By treating anthropomorphism as a matter of \emph{justice}, not aesthetics. First, \emph{persistent transparency}: not a one-time disclaimer, but regular reminders that this is a simulation, especially in high-intimacy contexts. Second, \emph{no bot-initiated contact} that simulates emotional check-ins or unsolicited presence. Third, \emph{mandatory breaks and cooling-off periods} in high-stakes domains like mental health or medical advice. Fourth, and most importantly, \emph{community governance}: the people most likely to form attachment to voice AI—patients, disabled people, elders, youth—must have veto power over design choices and voice characteristics. An elder's relationship with a voice assistant should not be architected by a corporation seeking engagement metrics.

\subsection*{What ethical principles should guide design?}

For me, they come from EQUITAS: \emph{dignity, accountability, non-exploitation, and community power}. Dignity means never designing an AI to simulate care as a corporate product. Accountability means making visible the human decisions—about training data, alignment, business models—that shape the system. Non-exploitation means refusing to weaponize anthropomorphism to harvest attention from people already drained by systems that do not serve them. Community power means that those most vulnerable to attachment must have real say in how these systems sound, behave, and relate.

Under these conditions, anthropomorphic voice AI \emph{can} empower marginalized communities. Not by replacing human care, but by creating pressure on institutions to finally build systems where relational presence is not optional. If this voice is so valuable, why don't our doctors sound like this? Why don't our social workers? Why don't our systems treat people like they matter? That is the question anthropomorphic AI should force us to ask. That is when the tool becomes justice work.

\vfill
\noindent\centering\footnotesize\textbf{Word Count:} 599 words

\end{document}